% =========================================================================
%  Figura: los cuatro casos de conductor hueco (como en las figuras del tema).
%  Contornos en polares (centro común):
%    cara exterior  r(t) = 1.50 + 0.16 cos 3t + 0.10 sen 2t
%    gaussiana      r(t) = 1.06 + 0.08 cos 3t + 0.05 sen 2t
%    cara interior  r(t) = 0.62 + 0.08 cos(2t + 40)
% =========================================================================
\newcommand{\cavExt}{plot[domain=0:360, samples=90, smooth cycle, variable=\t]
  ({(1.50+0.16*cos(3*\t)+0.10*sin(2*\t))*cos(\t)},{(1.50+0.16*cos(3*\t)+0.10*sin(2*\t))*sin(\t)})}
\newcommand{\cavInt}{plot[domain=0:360, samples=70, smooth cycle, variable=\t]
  ({(0.62+0.08*cos(2*\t+40))*cos(\t)},{(0.62+0.08*cos(2*\t+40))*sin(\t)})}
\newcommand{\cavGauss}{plot[domain=0:360, samples=80, smooth cycle, variable=\t]
  ({(1.06+0.08*cos(3*\t)+0.05*sin(2*\t))*cos(\t)},{(1.06+0.08*cos(3*\t)+0.05*sin(2*\t))*sin(\t)})}
% signos sobre la cara exterior (#1 = 1) o interior (#1 = 2), en el lado del metal
\newcommand{\signosExt}{\foreach \t in {0,30,...,330}{
  \node[signopos, font=\tiny\bfseries] at ({(1.36+0.16*cos(3*\t)+0.10*sin(2*\t))*cos(\t)},{(1.36+0.16*cos(3*\t)+0.10*sin(2*\t))*sin(\t)}) {$+$};}}
\newcommand{\signosInt}{\foreach \t in {0,36,...,324}{
  \node[signoneg, font=\tiny\bfseries] at ({(0.76+0.08*cos(2*\t+40))*cos(\t)},{(0.76+0.08*cos(2*\t+40))*sin(\t)}) {$-$};}}
\newcommand{\panelCavidad}[4]{% #1 título, #2 signos ext (0/1), #3 carga en cavidad (0/1), #4 rótulos
  \fill[metal!55!white, even odd rule] \cavExt \cavInt;
  \fill[pattern=north east lines, pattern color=azulcielo!70, even odd rule] \cavExt \cavInt;
  \draw[azulprof, line width=0.6pt] \cavExt;
  \draw[azulprof, line width=0.6pt] \cavInt;
  \draw[gauss] \cavGauss;
  \ifnum#2=1 \signosExt\fi
  \ifnum#3=1 \signosInt \node[qpos, minimum size=8pt] at (0.05,-0.02) {$+$};
             \node[font=\scriptsize, text=cargapos] at (0.05,0.3) {$q$};\fi
  \node[font=\scriptsize, text=azulprof, fill=white, inner sep=1pt, rounded corners=1pt] at (0.15,-1.25) {$\vE=\vc{0}$};
  \node[anchor=west, align=left, font=\footnotesize] at (1.95,0) {#4};
  \node[anchor=south, font=\small] at (1.15,1.85) {#1};}

\begin{figure}[t]
\centering
\begin{tikzpicture}[x=0.95cm, y=0.95cm]
  \begin{scope}
    \panelCavidad{(a) Neutro, cavidad vacía}{0}{0}
      {$\Qint=0$\\ $\Qext=0$\\[2pt] $\Ecav=\vc{0}$\\ $\Eext=\vc{0}$}
  \end{scope}
  \begin{scope}[xshift=8.2cm]
    \panelCavidad{(b) Carga $Q$, cavidad vacía}{1}{0}
      {$\Qint=0$\\ $\Qext=Q$\\[2pt] $\Ecav=\vc{0}$\\ $\Eext\neq\vc{0}$}
  \end{scope}
  \begin{scope}[yshift=-4.4cm]
    \panelCavidad{(c) Neutro, carga $q$ en la cavidad}{1}{1}
      {$\Qint=-q$\\ $\Qext=+q$\\[2pt] $\Ecav\neq\vc{0}$\\ $\Eext\neq\vc{0}$}
  \end{scope}
  \begin{scope}[xshift=8.2cm, yshift=-4.4cm]
    \panelCavidad{(d) Carga $Q$, carga $q$ en la cavidad}{1}{1}
      {$\Qint=-q$\\ $\Qext=Q+q$\\[2pt] $\Ecav\neq\vc{0}$\\ $\Eext\neq\vc{0}$}
  \end{scope}
  % leyenda común
  \draw[gauss] (2.6,-7.05) -- (3.4,-7.05);
  \node[anchor=west, font=\footnotesize] at (3.5,-7.05) {superficie de Gauss (cerrada) dentro del metal};
\end{tikzpicture}
\caption{Conductores huecos en equilibrio. La superficie de Gauss trazada dentro del metal, donde $\vE=\vc{0}$, encierra carga nula: la cara
interior lleva siempre la carga opuesta a la de la cavidad, y la exterior el resto de la carga del conductor. Una cavidad vacía no tiene campo.}
\label{fig:cavidades}
\end{figure}
